\documentclass[11pt]{article}
\usepackage[T1]{fontenc}
\usepackage[utf8]{inputenc}
\usepackage{lmodern}
\usepackage{microtype}
\usepackage[margin=1in]{geometry}
\usepackage{booktabs,tabularx,array}
\usepackage[hidelinks]{hyperref}
\usepackage[dvipsnames]{xcolor}
\usepackage{enumitem}
\usepackage{tcolorbox}
\setlength{\parskip}{0.55em}
\setlength{\parindent}{0pt}
\renewcommand{\arraystretch}{1.15}
\title{Phase 6 report\\\large Hostile export gates passed for the kernel-first same-domain Godel manuscript}
\author{OpenAI}
\date{\today}
\begin{document}
\maketitle

\begin{abstract}
Phase~6 executes the rewrite plan's hostile-export battery. The primary manuscript is audited against six explicit risks: circularity, syntax compression, provisional-practice misclassification, scope drift, loss of Godel-specificity, and proof-ineconomy. The result of this phase is a gate-passed primary-manuscript seed that remains kernel-first. All downstream interface, quotient, and coherent-reasoning layers are still treated as forced consequences relative to that kernel, and the coherent-reasoning layer remains the decisive defeat of the hardest hidden-path attack.
\end{abstract}

\section{Phase-6 objective}
The rewrite plan's Phase~6 instruction is simple and strict: no draft counts as primary-manuscript ready until it passes the hostile export gates explicitly. In the present project that means checking whether a hostile reader can still reinterpret the paper as circular, as mere syntax relabeling, as a ban on ordinary mathematical revision, as a theorem about all calculi, as insufficiently Godel-specific, or as obscuring its shortest theorem path under audit burden.

\section{Kernel focus preserved}
Phase~6 does not change the architecture fixed in earlier phases. The manuscript remains \emph{kernel-first}. Appendix~A fixes the compact floor, and all downstream interface, quotient, and coherent-reasoning layers remain treated as forced consequences \emph{relative to that kernel}. The point of this phase is not to reopen architecture but to ensure that the manuscript now exports that architecture cleanly.

\begin{tcolorbox}[title=Hardest hostile attack closed]
The strongest remaining hostile move says that a same-domain Godel threat might matter only through intermediate derivation paths, proof-theoretic strength, or deferred trajectory structure without surfacing in standing classification or endpoint behavior. The coherent-reasoning layer blocks exactly that route. In the fixed regime, coherence is pointwise standing-preservation along legal same-scope trajectories; incoherent prefixes remain incoherent under admissible continuation; terminal standing fully certifies coherence; all coherence-relevant distinctions are tensor-determined rather than skin-dependent; and standing failure is irrecoverable. There is therefore no third hidden inferential locus left for a same-domain threat to inhabit.
\end{tcolorbox}

\section{Hostile export gate matrix}
\begin{center}
\begin{tabularx}{\textwidth}{>{\bfseries}p{0.22\textwidth} >{\raggedright\arraybackslash}p{0.15\textwidth} X}
\toprule
Gate & Status & Evidence in the revised seed \\
\midrule
Circularity gate & PASS after audit & Appendix~A still fixes the compact kernel floor before the body application begins. The body starts only after inferential-life exhaustion and same-domain threat requirements are fixed. \\
Syntax-compression gate & PASS after repair & The introduction now states explicitly that syntax, coding, and proof inscriptions matter only through same-domain readback into standing classification or coherent trajectory structure. Section~4 adds an explicit syntax-compression remark and a bare-syntax corollary. \\
Provisional-practice gate & PASS after repair & The limits section now classifies seminar sketches, drafts, repaired proofs, and provisional codings as provisional practice until standing-bearing uptake, or else as fresh acts or explicit resets. \\
Scope gate & PASS after audit & Abstract, introduction, body definitions, and limits all describe the same target: same-domain governance-relevant threat architectures inside a fixed kernel-relative regime. \\
Godel-specificity gate & PASS after audit & Sections~3--6 remain application-directed: threat requirements, code-role exhaustion, no same-domain diagonalization, rescue-route exhaustion, and the terminal theorem. \\
Proof-economy gate & PASS after repair & The roadmap now states the exact eight-result visible body path and keeps the audit reservoir outside the primary reader path. \\
\bottomrule
\end{tabularx}
\end{center}

\section{Concrete repairs made}
Phase~6 makes four targeted repairs.
\begin{enumerate}[leftmargin=2.2em]
    \item The introduction now explicitly states that the target is not bare derivation-trace classification.
    \item Section~4 now includes an explicit syntax-compression remark and a bare-syntax corollary.
    \item The limits section now classifies provisional mathematical practice and repaired proofs without weakening the no-repair theorem.
    \item The roadmap now names the exact eight-result visible body path, making proof economy explicit rather than tacit.
\end{enumerate}

\section{Why the manuscript now passes the six gates}
\subsection*{Circularity}
The application no longer appears to define its own regime. Appendix~A fixes the kernel floor. Appendices~B--E record forced consequences relative to that floor. The body then begins with the inferential-life bridge and same-domain threat requirements. The Godel application is therefore plainly downstream.

\subsection*{Syntax compression}
A hostile formalist can no longer honestly compress the thesis to ``syntax relabeled as standing.'' The revised introduction and Section~4 now state directly that a coding difference with no governance-relevant readback is bookkeeping, not a same-domain closure threat.

\subsection*{Provisional practice}
This was the most likely remaining misread. The revised limits section now states cleanly that seminar sketches, drafts, repaired proofs, and provisional codings are not forbidden. They simply do not count as standing-bearing same-domain acts until taken up as such. Once they are so taken up, later corrections are fresh acts or explicit resets, not same-act rescue.

\subsection*{Scope consistency}
The project now maintains one target description throughout: same-domain governance-relevant Godel-threat architecture inside a fixed kernel-relative regime. No page-1 sentence suggests a theorem about all reasoning practices or all formal calculi.

\subsection*{Godel-specificity}
The manuscript remains visibly a Godel-response paper because Sections~3--6 still track the exact application path: threat requirements, code-role exhaustion, diagonal-route failure, rescue-route exhaustion, terminal non-instantiability.

\subsection*{Proof economy}
The visible body path is now explicit and finite. The report does not deny the theorem reservoir. It shows only that the reservoir is no longer the reader's required path through the primary manuscript.

\section{Deliverables in this phase}
The Overleaf project produced in Phase~6 includes:
\begin{itemize}
    \item \texttt{phase6\_seed.tex}: the hostile-export-gated primary-manuscript seed;
    \item the revised body section files under \texttt{seed/sections/};
    \item the compact appendix stack under \texttt{seed/appendices/};
    \item \texttt{data/phase6\_gate\_matrix.csv} and \texttt{data/phase6\_gate\_matrix.json};
    \item the continuing audit companion and frozen audit reserve; and
    \item this report as \texttt{main.tex} for immediate Overleaf compilation.
\end{itemize}

\section{Exit status}
Phase~6 is complete. All six hostile export gates now pass at the level of the primary-manuscript seed. The manuscript remains kernel-first, all downstream layers remain treated as forced consequences relative to that kernel, and the coherent-reasoning layer is explicitly positioned as the defeat of the hardest hidden-path attack.

\end{document}
